\subsection{平衡外置的 Schur 最优性、临界审计阶梯与二原子 Mellin 平台律}\label{subsec:conclusion-balanced-externalization-harmonic-staircase-mellin-plateau}
沿用命题 \ref{prop:fold-oracle-variational-projection-characterization}、定理 \ref{thm:fold-bin-two-state-asymptotic}、定理 \ref{thm:conclusion-foldbin-oracle-critical-profile}、定理 \ref{thm:conclusion-binfold-mellin-two-step-law} 与定理 \ref{thm:conclusion-binfold-uniform-output-twoatom-choquet} 的记号。前述链条已经分别把预算约束下的预言机容量、bin-fold 纤维规模的两态一致渐近、critical-scale 容量曲线以及二原子 Mellin--Choquet 机制固定为闭式。若再把纤维内外置的离散平衡原理与容量剖面的 Mellin 表达并读，则还可进一步抽出如下几条刚性后果。

\begin{theorem}[纤维内平衡外置是全部 Schur-convex 审计泛函的共同极小者]\label{thm:conclusion-balanced-externalization-schur-optimality}
\leanverified{paper\_conclusion\_balanced\_externalization\_schur\_optimality}
固定分辨率 \(m\) 与外置位数 \(B\)，并记
\[
M:=2^B,\qquad
\Omega_m:=\{0,1\}^m,
\qquad
d_m(x):=\bigl|(\mathrm{Fold}_m)^{-1}(x)\bigr|.
\]
对任意外置映射
\[
R:\Omega_m\to\{0,1\}^B,
\]
定义每条纤维 \(x\in X_m\) 上的占用向量
\[
n_x(R):=\bigl(n_{x,r}(R)\bigr)_{r\in\{0,1\}^B},
\qquad
n_{x,r}(R):=\#\{\omega\in\Omega_m:\ \mathrm{Fold}_m(\omega)=x,\ R(\omega)=r\}.
\]
若对整数 \(d\ge 0\) 写
\[
d=aM+r,\qquad 0\le r<M,
\]
并记 \(\mathbf b(d,M)\) 为由 \(r\) 个 \(a+1\) 与 \(M-r\) 个 \(a\) 组成的平衡向量。若
\[
\Psi:\RR_{\ge 0}^{M}\to\RR
\]
是任意对称 Schur-convex 泛函，则
都有
\[
\sum_{x\in X_m}\Psi\!\bigl(n_x(R)\bigr)
\ge
\sum_{x\in X_m}\Psi\!\bigl(\mathbf b(d_m(x),M)\bigr).
\]
并且存在某个外置映射 \(R_{\mathrm{bal}}\) 同时在每条纤维上达到该下界。
\end{theorem}
\begin{proof}
固定一条纤维 \(x\in X_m\)，并记 \(d:=d_m(x)\)。在所有满足
\[
(n_1,\dots,n_M)\in \ZZ_{\ge 0}^M,
\qquad
\sum_{j=1}^{M}n_j=d
\]
的占用向量中，\(\mathbf b(d,M)\) 是 majorization 意义下的极小元；等价地，任意这类向量都 majorizes \(\mathbf b(d,M)\)。因此
\[
(n_1,\dots,n_M)\succcurlyeq \mathbf b(d,M).
\]
由 Schur-convex 性，进而有
\[
\Psi(n_1,\dots,n_M)\ge \Psi\!\bigl(\mathbf b(d,M)\bigr).
\]
对每个 \(x\in X_m\) 独立应用此结论，即得逐纤维下界；再对全部纤维求和便得到所示全局不等式。

另一方面，不同纤维之间不存在耦合。故可在每条纤维内部独立选取一个把 \(d_m(x)\) 个元素尽可能均匀分配到 \(M\) 个桶中的映射；将这些局部选择并合，便得到全局外置映射 \(R_{\mathrm{bal}}\)，并使每条纤维都恰好实现 \(\mathbf b(d_m(x),M)\)。故下界全局可达。证毕。
\end{proof}

\begin{corollary}[固定比特预算下最优 Rényi 型碰撞能量的精确闭式]\label{cor:conclusion-balanced-externalization-optimal-renyi-energy}
\leanverified{paper\_conclusion\_balanced\_externalization\_optimal\_renyi\_energy}
对任意实数 \(q\ge 1\)，定义
\[
\mathcal J_{q,m}^{\mathrm{opt}}(B)
:=
\min_{R:\Omega_m\to\{0,1\}^B}
\sum_{x\in X_m}\sum_{r\in\{0,1\}^B} n_{x,r}(R)^q.
\]
若对每个 \(x\in X_m\) 写
\[
d_m(x)=a_x\,2^B+r_x,\qquad 0\le r_x<2^B,
\]
则有精确公式
\[
\mathcal J_{q,m}^{\mathrm{opt}}(B)
=
\sum_{x\in X_m}
\Bigl[
r_x(a_x+1)^q+(2^B-r_x)a_x^q
\Bigr].
\]
特别地，最优外置机制与 \(q\) 无关；同一平衡外置同时最小化全部 \(q\ge 1\) 的 Rényi 型碰撞能量。
\end{corollary}
\begin{proof}
在定理 \ref{thm:conclusion-balanced-externalization-schur-optimality} 中取
\[
\Psi(v_1,\dots,v_{2^B})=\sum_{j=1}^{2^B}v_j^q
\]
即可。由于 \(q\ge 1\) 时 \(t\mapsto t^q\) 为凸函数，故上述 \(\Psi\) 是对称 Schur-convex 的。再把平衡向量 \(\mathbf b(d_m(x),2^B)\) 的分量写成 \(r_x\) 个 \(a_x+1\) 与 \(2^B-r_x\) 个 \(a_x\)，便得到所示闭式。证毕。
\end{proof}

\begin{theorem}[bin-fold 临界预算下最优 \(q\)-能量密度的极限相图]\label{thm:conclusion-binfold-critical-budget-renyi-energy}
\leanverified{paper\_conclusion\_binfold\_critical\_budget\_renyi\_energy}
考虑 bin-fold 体系 \(\mathrm{Fold}_m^{\mathrm{bin}}\)，并记
\[
A_m:=\Bigl(\frac{2}{\varphi}\Bigr)^m,
\qquad
\mathcal E_{q,m}(B):=2^{-m}\mathcal J_{q,m}^{\mathrm{opt}}(B).
\]
设整数预算序列 \((B_m)_{m\ge 1}\) 满足
\[
2^{B_m}=\tau A_m(1+o(1)),
\qquad
\tau>0.
\]
再定义
\[
\Theta_q(\rho):=
\bigl(1-\{\rho\}\bigr)\lfloor\rho\rfloor^q
+
\{\rho\}\,(\lfloor\rho\rfloor+1)^q,
\qquad
\{\rho\}:=\rho-\lfloor\rho\rfloor.
\]
则当 \(m\to\infty\) 时，
\[
\mathcal E_{q,m}(B_m)\longrightarrow \mathcal E_q(\tau),
\]
其中
\[
\mathcal E_q(\tau)
=
\frac{\varphi}{\sqrt5}\,\tau\,\Theta_q\!\Bigl(\frac1\tau\Bigr)
+
\frac1{\sqrt5}\,\tau\,\Theta_q\!\Bigl(\frac{\varphi^{-1}}{\tau}\Bigr).
\]
更具体地，对每个 \(c\in\{1,\varphi^{-1}\}\)，函数
\[
G_{q,c}(\tau):=\tau\,\Theta_q(c/\tau)
\]
在每个区间
\[
\tau\in\Bigl[\frac{c}{n+1},\frac{c}{n}\Bigr]
\qquad (n\ge 1)
\]
上都是仿射函数，并满足
\[
G_{q,c}(\tau)
=
c\bigl((n+1)^q-n^q\bigr)
+
\tau\bigl((n+1)n^q-n(n+1)^q\bigr).
\]
当 \(\tau\ge c\) 时，则有
\[
G_{q,c}(\tau)=c.
\]
\end{theorem}
\begin{proof}
先记单纤维最优 \(q\)-能量为
\[
\psi_q(d,M):=
\min_{\substack{(n_1,\dots,n_M)\in\ZZ_{\ge 0}^M\\ n_1+\cdots+n_M=d}}
\sum_{j=1}^{M}n_j^q.
\]
若 \(d=aM+r\) 且 \(0\le r<M\)，则由推论 \ref{cor:conclusion-balanced-externalization-optimal-renyi-energy}
\[
\psi_q(d,M)=r(a+1)^q+(M-r)a^q.
\]
等价地，写 \(\rho=d/M=a+r/M\)，便有精确恒等式
\[
\psi_q(d,M)=M\,\Theta_q(\rho)=M\,\Theta_q(d/M).
\]

现在对 bin-fold 取
\[
A_0^{(m)}:=\{w\in X_m:\ w_m=0\},
\qquad
A_1^{(m)}:=\{w\in X_m:\ w_m=1\}.
\]
由定理 \ref{thm:fold-bin-two-state-asymptotic} 的一致展开，
\[
d_m^{\mathrm{bin}}(w)=A_m+o(A_m)
\quad (w\in A_0^{(m)}),
\qquad
d_m^{\mathrm{bin}}(w)=\varphi^{-1}A_m+o(A_m)
\quad (w\in A_1^{(m)}),
\]
且误差在各自扇区上一致。又
\[
|A_0^{(m)}|=F_{m+1},
\qquad
|A_1^{(m)}|=F_m.
\]
记 \(M_m:=2^{B_m}=\tau A_m(1+o(1))\)。则对任意固定 \(c>0\)，若 \(d_m=cA_m+o(A_m)\)，便有
\[
\frac{\psi_q(d_m,M_m)}{A_m}
=
\frac{M_m}{A_m}\,\Theta_q(d_m/M_m)
\longrightarrow
\tau\,\Theta_q(c/\tau),
\]
其中用到 \(\Theta_q\) 在 \([0,\infty)\) 上连续。

于是
\[
\mathcal E_{q,m}(B_m)
=
2^{-m}\sum_{w\in A_0^{(m)}}\psi_q(d_m(w),M_m)
+
2^{-m}\sum_{w\in A_1^{(m)}}\psi_q(d_m(w),M_m).
\]
再由
\[
\frac{F_{m+1}A_m}{2^m}\longrightarrow \frac{\varphi}{\sqrt5},
\qquad
\frac{F_mA_m}{2^m}\longrightarrow \frac1{\sqrt5},
\]
即可得到
\[
\mathcal E_{q,m}(B_m)\longrightarrow
\frac{\varphi}{\sqrt5}\,\tau\,\Theta_q\!\Bigl(\frac1\tau\Bigr)
+
\frac1{\sqrt5}\,\tau\,\Theta_q\!\Bigl(\frac{\varphi^{-1}}{\tau}\Bigr).
\]

若
\[
\tau\in\Bigl[\frac{c}{n+1},\frac{c}{n}\Bigr],
\]
则 \(c/\tau\in[n,n+1]\)。写
\[
\frac{c}{\tau}=n+\theta,
\qquad
0\le \theta\le 1.
\]
于是
\[
\Theta_q(c/\tau)=(1-\theta)n^q+\theta(n+1)^q,
\qquad
\theta=\frac{c}{\tau}-n.
\]
代回并整理即得
\[
G_{q,c}(\tau)
=
c\bigl((n+1)^q-n^q\bigr)
+
\tau\bigl((n+1)n^q-n(n+1)^q\bigr).
\]
若 \(\tau\ge c\)，则 \(0\le c/\tau\le 1\)，从而 \(\Theta_q(c/\tau)=c/\tau\)，于是 \(G_{q,c}(\tau)=c\)。证毕。
\end{proof}

\begin{corollary}[高阶审计曲线的无穷谐波折点结构]\label{cor:conclusion-binfold-harmonic-staircase-kinks}
若 \(q>1\)，则极限曲线 \(\mathcal E_q(\tau)\) 的全部不可微点恰为
\[
\Sigma_q=
\Bigl\{\frac1n:\ n\ge 1\Bigr\}
\cup
\Bigl\{\frac{\varphi^{-1}}{n}:\ n\ge 1\Bigr\}.
\]
该集合在 \(0\) 处聚积，且除此之外无其他聚积点。因而 bin-fold 的临界外置几何并非有限折点图景，而是一条向 \(0\) 积聚的无穷谐波审计阶梯。
\end{corollary}
\begin{proof}
由定理 \ref{thm:conclusion-binfold-critical-budget-renyi-energy}，
\[
\mathcal E_q(\tau)
=
\frac{\varphi}{\sqrt5}\,G_{q,1}(\tau)
+
\frac1{\sqrt5}\,G_{q,\varphi^{-1}}(\tau).
\]
对固定 \(c>0\)，\(G_{q,c}\) 只可能在
\[
\tau=\frac{c}{n}
\qquad (n\ge 1)
\]
处发生折点，因为这些点正对应于 \(\lfloor c/\tau\rfloor\) 的跳变。并且在区间
\[
\Bigl[\frac{c}{n+1},\frac{c}{n}\Bigr]
\]
上的斜率等于
\[
\sigma_{q,n}:=(n+1)n^q-n(n+1)^q.
\]
于是
\[
\sigma_{q,n}-\sigma_{q,n-1}
=
-n\Bigl((n+1)^q-2n^q+(n-1)^q\Bigr)<0,
\]
其中严格不等式来自 \(t\mapsto t^q\) 在 \(q>1\) 时的严格凸性。故每个 \(\tau=c/n\) 都是真正折点。

取 \(c=1\) 与 \(c=\varphi^{-1}\)，便得全部折点包含于
\[
\Bigl\{\frac1n:\ n\ge 1\Bigr\}
\cup
\Bigl\{\frac{\varphi^{-1}}{n}:\ n\ge 1\Bigr\}.
\]
由于 \(\varphi^{-1}\notin\QQ\)，两条谐波格彼此不重合，故总折点集恰为其并。显然其唯一聚积点为 \(0\)。证毕。
\end{proof}
\leanverified{paper\_conclusion\_binfold\_harmonic\_staircase\_kinks}

\begin{theorem}[负矩 Mellin 指纹对二原子容量律的唯一恢复]\label{thm:conclusion-negative-mellin-two-atom-uniqueness}
\leanverified{paper\_conclusion\_negative\_mellin\_two\_atom\_uniqueness}
设 \(\nu\) 为支撑有界的 Borel 概率测度，并定义其容量剖面
\[
m_\nu(\tau):=\int \min\{y,\tau\}\,\dd\nu(y),
\qquad
\tau\ge 0.
\]
若对每个 \(s>0\) 都有
\[
s(s+1)\int_0^\infty \tau^{-s-2}\bigl(\tau-m_\nu(\tau)\bigr)\,\dd\tau
=
\varphi^{-1}+\varphi^{\,s-2},
\]
则必有
\[
\nu=\varphi^{-2}\delta_{\varphi^{-1}}+\varphi^{-1}\delta_1.
\]
换言之，二原子 scaled multiplicity law 已经被其负矩 Mellin 指纹单独唯一刻画。
\end{theorem}
\begin{proof}
对任意固定 \(y>0\)，有
\[
\tau-\min\{y,\tau\}=
\begin{cases}
0, & 0<\tau\le y,\\[2pt]
\tau-y, & \tau>y.
\end{cases}
\]
故直接计算得到
\[
s(s+1)\int_0^\infty \tau^{-s-2}\bigl(\tau-\min\{y,\tau\}\bigr)\,\dd\tau
=
y^{-s}.
\]
对 \(\nu\) 积分并交换积分次序，即得
\[
s(s+1)\int_0^\infty \tau^{-s-2}\bigl(\tau-m_\nu(\tau)\bigr)\,\dd\tau
=
\int y^{-s}\,\dd\nu(y).
\]
于是题设等价于
\[
\int y^{-s}\,\dd\nu(y)
=
\varphi^{-1}+\varphi^{-2}(\varphi^{-1})^{-s}
\qquad (s>0).
\]
记
\[
\mu:=\nu-\bigl(\varphi^{-2}\delta_{\varphi^{-1}}+\varphi^{-1}\delta_1\bigr).
\]
则 \(\mu\) 是紧支撑有限符号测度，并满足
\[
\int y^{-s}\,\dd\mu(y)=0
\qquad (\forall\,s>0).
\]
将 \(\mu\) 通过变量代换 \(y=e^t\) 推前为一支撑紧的有限符号测度 \(\eta\)，便得到
\[
\int e^{-st}\,\dd\eta(t)=0
\qquad (\forall\,s>0).
\]
由紧支撑有限测度的 Laplace 变换唯一性，\(\eta=0\)，故 \(\mu=0\)。因此
\[
\nu=\varphi^{-2}\delta_{\varphi^{-1}}+\varphi^{-1}\delta_1.
\]
证毕。
\end{proof}

\begin{theorem}[completion 的二值分支并不自足；容量斜率比是唯一选枝器]\label{thm:conclusion-balanced-twoatom-completion-capacity-branch-selection}
\leanverified{paper\_conclusion\_balanced\_twoatom\_completion\_capacity\_branch\_selection}
对任意 \(p\in(0,1)\)，定义两原子概率测度
\[
\nu_{\varphi,p}=p\,\delta_1+(1-p)\,\delta_{\varphi^{-1}},
\qquad
M_{\varphi,p}(z)=p+(1-p)\varphi^z.
\]
则存在常数 \(\kappa>0\) 使
\[
\kappa\,M_{\varphi,p}(2s)M_{\varphi,p}(1-2s)
=
\bigl(1+\varphi^{2s}\bigr)\bigl(1+\varphi^{1-2s}\bigr)
\qquad (\forall s\in\CC)
\]
当且仅当
\[
(p,\kappa)=\left(\frac12,4\right)
\qquad\text{或}\qquad
(p,\kappa)=\left(\varphi^{-1},\varphi^3\right).
\]
也就是说，completion 自对偶形式本身只把 projective 两原子律钉到一个二值分支。

但若再要求由 \(\nu_{\varphi,p}\) 的 \(\min\)-变换所给出的双折点容量曲线，与项目临界容量曲线在两段线性区间上的斜率比完全一致，则唯一允许的参数立即坍缩为
\[
p=\varphi^{-1},
\qquad
\nu_{\varphi,p}=\varphi^{-1}\delta_1+\varphi^{-2}\delta_{\varphi^{-1}}.
\]
因此，completion 数据本身并不能唯一恢复项目尺度律；真正完成唯一选枝的是容量侧的折点几何。
\end{theorem}
\begin{proof}
正文已经完成了这条二分链的全部计算。把
\[
\kappa\,M_{\varphi,p}(2s)M_{\varphi,p}(1-2s)
\]
展开后，与
\[
\bigl(1+\varphi^{2s}\bigr)\bigl(1+\varphi^{1-2s}\bigr)
\]
逐项比较，正文得到唯一可能的参数对恰为
\[
\left(\frac12,4\right)
\qquad\text{或}\qquad
\left(\varphi^{-1},\varphi^3\right).
\]
因此，仅从 completion 自对偶因子出发，确实存在二值分支歧义。

另一方面，正文又把 \(\nu_{\varphi,p}\) 的 \(\min\)-变换容量曲线写成一条在 \(s=\varphi^{-1}\) 处发生折点的双线性函数，并证明其左右斜率比恰等于 \(p\)。而项目真实的临界容量曲线已经被另一条主链固定，其斜率比正是 \(\varphi^{-1}\)。故若要求 completion 与容量相图同时相容，则只能有
\[
p=\varphi^{-1}.
\]
于是二值歧义被容量几何唯一消去。证毕。
\end{proof}


\leanverified{paper\_conclusion\_supercritical\_plateau\_subcritical\_power\_law}
\begin{theorem}[超临界精确平台与亚临界幂律爆破]\label{thm:conclusion-supercritical-plateau-subcritical-power-law}
若 \(q>1\)，则定理 \ref{thm:conclusion-binfold-critical-budget-renyi-energy} 的极限曲线满足
\[
\mathcal E_q(\tau)\equiv 1
\qquad (\tau\ge 1).
\]
也就是说，一旦外置预算跨过最大原子尺度，最优高阶碰撞密度便出现严格平台。另一方面，当 \(\tau\downarrow 0\) 时，
\[
\mathcal E_q(\tau)
=
\frac{\varphi+\varphi^{-q}}{\sqrt5}\,\tau^{1-q}
+
O(\tau^{2-q}).
\]
特别地，在二阶碰撞情形 \(q=2\) 下，
\[
\mathcal E_2(\tau)
=
\frac{2}{\sqrt5}\,\tau^{-1}+O(1),
\qquad
\tau\downarrow 0.
\]
\end{theorem}
\begin{proof}
若 \(\tau\ge 1\)，则
\[
\frac1\tau\le 1,
\qquad
\frac{\varphi^{-1}}{\tau}\le \varphi^{-1}<1.
\]
当 \(0\le \rho\le 1\) 时，\(\Theta_q(\rho)=\rho\)。因此
\[
\mathcal E_q(\tau)
=
\frac{\varphi}{\sqrt5}\,\tau\cdot \frac1\tau
+
\frac1{\sqrt5}\,\tau\cdot \frac{\varphi^{-1}}{\tau}
=
\frac{\varphi+\varphi^{-1}}{\sqrt5}
=1.
\]

若 \(\tau\downarrow 0\)，则 \(\rho\to\infty\) 时有
\[
\Theta_q(\rho)=\rho^q+O(\rho^{q-1}),
\]
因为写 \(\rho=n+\theta\) 且 \(0\le \theta<1\) 时，
\[
\Theta_q(\rho)=(1-\theta)n^q+\theta(n+1)^q
=
\rho^q+O(\rho^{q-1}).
\]
于是
\[
\tau\,\Theta_q(1/\tau)=\tau^{1-q}+O(\tau^{2-q}),
\]
\[
\tau\,\Theta_q(\varphi^{-1}/\tau)=\varphi^{-q}\tau^{1-q}+O(\tau^{2-q}).
\]
代入定理 \ref{thm:conclusion-binfold-critical-budget-renyi-energy} 即得
\[
\mathcal E_q(\tau)
=
\frac{\varphi+\varphi^{-q}}{\sqrt5}\,\tau^{1-q}
+
O(\tau^{2-q}).
\]
对 \(q=2\)，再用恒等式
\[
\varphi+\varphi^{-2}=2
\]
便得到
\[
\mathcal E_2(\tau)=\frac{2}{\sqrt5}\tau^{-1}+O(1).
\]
证毕。
\end{proof}

\paragraph{统一读法}
这组结果把 bin-fold 外置预算几何压缩为一条比连续热力学图像更刚性的离散相图：在固定窗口上，纤维内平衡外置同时最小化全部凸审计泛函与全部 \(q\)-碰撞能量；在临界尺度上，这一最优机制收敛为由两条谐波格
\[
\Bigl\{\frac1n\Bigr\}_{n\ge 1},
\qquad
\Bigl\{\frac{\varphi^{-1}}{n}\Bigr\}_{n\ge 1}
\]
支配的无穷折点阶梯；在 Mellin 侧，整个容量剖面又被一支二原子测度完全唯一决定，并强制导出超临界平台与亚临界幂律爆破。进一步地，completion 自对偶因子虽然单独只给出一个二值分支，但一旦把容量相图的折点斜率比并入同一审计框架，分支立刻唯一坍缩到
\[
\varphi^{-1}\delta_1+\varphi^{-2}\delta_{\varphi^{-1}}.
\]
于是，budgeted externalization、critical-scale 容量、Mellin 指纹与 quarter-line completion 四者不再只是平行出现的统计像，而被同一两原子几何唯一锁定。

\endinput
